\documentclass[11pt]{article}

\usepackage[utf8]{inputenc}
\usepackage[T1]{fontenc}
\usepackage[a4paper,margin=2.5cm]{geometry}
\usepackage{amsmath}
\usepackage{natbib}
\usepackage{hyperref}
\usepackage{filecontents}

%==========================================================================
%  Embedded BibTeX entries: writes the file `intro_additions.bib` on the
%  first compile.  Compile sequence:   pdflatex -> bibtex -> pdflatex x2
%==========================================================================
\begin{filecontents*}[overwrite]{intro_additions.bib}

%---------- Most important omission ----------

@misc{yozova2025data,
  author       = {Yozova, Ilina and McJames, Nathan and Manolopoulou, Ioanna},
  title        = {Data fusion under weak identifiability in heterogeneous treatment effect modelling},
  year         = {2025},
  note         = {O'Bayes Conference on Objective {B}ayes Methodology, Athens, Greece, June 2025. Conference presentation}
}

%---------- Already in working-paper references, candidates for Para 1 ----------

@article{chengcai2021adaptive,
  author       = {Cheng, Dan and Cai, Tianxi},
  title        = {Adaptive combination of randomized and observational data},
  journal      = {arXiv preprint arXiv:2111.15012},
  year         = {2021}
}

@article{shyr2025multi,
  author       = {Shyr, Cathy and Ren, Bo and Patil, Prasad and Parmigiani, Giovanni},
  title        = {Multi-study {R}-learner for estimating heterogeneous treatment effects across studies using statistical machine learning},
  journal      = {Biostatistics},
  volume       = {26},
  number       = {1},
  pages        = {kxaf040},
  year         = {2025}
}

@article{yangding2020combining,
  author       = {Yang, Shu and Ding, Peng},
  title        = {Combining multiple observational data sources to estimate causal effects},
  journal      = {Journal of the American Statistical Association},
  volume       = {115},
  number       = {531},
  pages        = {1540--1554},
  year         = {2020},
  doi          = {10.1080/01621459.2019.1609973}
}

@article{vanderlaan2024adaptive,
  author       = {van der Laan, Mark J. and Qiu, Shuangning and van der Laan, Lisha},
  title        = {Adaptive-{TMLE} for the average treatment effect based on randomized controlled trial augmented with real-world data},
  journal      = {arXiv preprint arXiv:2405.07186},
  year         = {2024}
}

@article{lintarpevans2024fusion,
  author       = {Lin, Xi and Tarp, Jens Magelund and Evans, Robin J.},
  title        = {Data fusion for efficiency gain in {ATE} estimation: A practical review with simulations},
  journal      = {arXiv preprint arXiv:2407.01186},
  year         = {2024},
  doi          = {10.48550/arXiv.2407.01186}
}

%---------- Already in working-paper references, candidates for Para 2 ----------

@article{muller2023bayesian,
  author       = {M{\"u}ller, Peter and Chandra, Noirrit Kiran and Sarkar, Abhra},
  title        = {Bayesian approaches to include real-world data in clinical studies},
  journal      = {Philosophical Transactions of the Royal Society A: Mathematical, Physical and Engineering Sciences},
  volume       = {381},
  number       = {2247},
  pages        = {20220158},
  year         = {2023},
  doi          = {10.1098/rsta.2022.0158}
}

@article{murray2015combining,
  author       = {Murray, Timothy A. and Hobbs, Brian P. and Carlin, Bradley P.},
  title        = {Combining nonexchangeable functional or survival data sources in oncology using generalized mixture commensurate priors},
  journal      = {The Annals of Applied Statistics},
  volume       = {9},
  number       = {3},
  pages        = {1549--1570},
  year         = {2015}
}

@article{ohigashi2022horseshoe,
  author       = {Ohigashi, Tetsuhisa and Ochi, Hisato and Aoki, Shintaro and Hamada, Chikuma},
  title        = {Using horseshoe prior for incorporating multiple historical control data in randomized controlled trials},
  journal      = {Statistical Methods in Medical Research},
  volume       = {31},
  number       = {7},
  pages        = {1392--1404},
  year         = {2022}
}

@article{ohigashi2025potential,
  author       = {Ohigashi, Tetsuhisa and Ochi, Hisato and Aoki, Shintaro and Hamada, Chikuma},
  title        = {Potential bias models with {Bayesian} shrinkage priors for dynamic borrowing of multiple historical control data},
  journal      = {Pharmaceutical Statistics},
  volume       = {24},
  number       = {2},
  pages        = {e2453},
  year         = {2025}
}

%---------- Already in methodology references, candidate for Para 2 promotion ----------

@article{henderson2020individualized,
  author       = {Henderson, Nicholas C. and Louis, Thomas A. and Rosner, Gary L. and Varadhan, Ravi},
  title        = {Individualized treatment effects with censored data via fully nonparametric {Bayesian} accelerated failure time models},
  journal      = {Biostatistics},
  volume       = {21},
  number       = {1},
  pages        = {50--68},
  year         = {2020},
  doi          = {10.1093/biostatistics/kxx075}
}

@article{hahn2020bayesian,
  author       = {Hahn, P. Richard and Murray, Jared S. and Carvalho, Carlos M.},
  title        = {Bayesian regression tree models for causal inference: Regularization, confounding, and heterogeneous effects},
  journal      = {Bayesian Analysis},
  volume       = {15},
  number       = {3},
  pages        = {965--1056},
  year         = {2020}
}

%---------- Foundational context already referenced elsewhere ----------

@article{stuart2011propensity,
  author       = {Stuart, Elizabeth A. and Cole, Stephen R. and Bradshaw, Catherine P. and Leaf, Philip J.},
  title        = {The use of propensity scores to assess the generalizability of results from randomized trials},
  journal      = {Journal of the Royal Statistical Society: Series A (Statistics in Society)},
  volume       = {174},
  number       = {2},
  pages        = {369--386},
  year         = {2011}
}

@article{dahabreh2019generalizing,
  author       = {Dahabreh, Issa J. and Robertson, Sarah E. and Tchetgen, Eric J. and Stuart, Elizabeth A. and Hern{\'a}n, Miguel A.},
  title        = {Generalizing causal inferences from individuals in randomized trials to all trial-eligible individuals},
  journal      = {Biometrics},
  volume       = {75},
  number       = {2},
  pages        = {685--694},
  year         = {2019}
}

@article{degtiar2023review,
  author       = {Degtiar, Irina and Rose, Sherri},
  title        = {A review of generalizability and transportability},
  journal      = {Annual Review of Statistics and Its Application},
  volume       = {10},
  number       = {1},
  pages        = {501--524},
  year         = {2023}
}

@article{antonelli2019high,
  author       = {Antonelli, Joseph and Parmigiani, Giovanni and Dominici, Francesca},
  title        = {High-dimensional confounding adjustment using continuous spike and slab priors},
  journal      = {Bayesian Analysis},
  volume       = {14},
  number       = {3},
  pages        = {805--828},
  year         = {2019},
  doi          = {10.1214/18-BA1138}
}

%---------- Not currently in references; verified citation details ----------

@article{hu2021estimating,
  author       = {Hu, Liangyuan and Ji, Jiayi and Li, Fan},
  title        = {Estimating heterogeneous survival treatment effect in observational data using machine learning},
  journal      = {Statistics in Medicine},
  volume       = {40},
  number       = {21},
  pages        = {4691--4713},
  year         = {2021},
  doi          = {10.1002/sim.9090}
}

@article{zhao2023efficient,
  author       = {Zhao, Pan and Josse, Julie and Yang, Shu},
  title        = {Efficient and robust transfer learning of optimal individualized treatment regimes with right-censored survival data},
  journal      = {Journal of Machine Learning Research},
  volume       = {26},
  year         = {2025},
  note         = {arXiv:2301.05491}
}

@article{colnet2024causal,
  author       = {Colnet, B{\'e}n{\'e}dicte and Mayer, Imke and Chen, Guanhua and Dieng, Awa and Li, Ruohong and Varoquaux, Ga{\"e}l and Vert, Jean-Philippe and Josse, Julie and Yang, Shu},
  title        = {Causal inference methods for combining randomized trials and observational studies: A review},
  journal      = {Statistical Science},
  year         = {2024},
  note         = {arXiv:2011.08047}
}

@article{bareinboim2016causal,
  author       = {Bareinboim, Elias and Pearl, Judea},
  title        = {Causal inference and the data-fusion problem},
  journal      = {Proceedings of the National Academy of Sciences},
  volume       = {113},
  number       = {27},
  pages        = {7345--7352},
  year         = {2016}
}

%---------- Suggested but BibTeX details to confirm ----------

@article{bi2023pamhc,
  author       = {Bi, Dehua and Zhou, Tianjian and Zhong, Wei and Ji, Yuan},
  title        = {{PAM-HC}: A {B}ayesian nonparametric construction of hybrid control for randomized clinical trials using external data},
  journal      = {arXiv preprint arXiv:2310.20087},
  year         = {2023}
}

@article{leeyangwang2022doubly,
  author       = {Lee, Dasom and Yang, Shu and Wang, Xiaofei},
  title        = {Doubly robust estimators for generalizing treatment effects on survival outcomes from randomized controlled trials to a target population},
  journal      = {Journal of Causal Inference},
  year         = {2022},
  note         = {Citation details to be verified}
}

@article{alt2023leap,
  author       = {Alt, Ethan M. and others},
  title        = {{LEAP}: The latent exchangeability prior for borrowing information from historical data},
  year         = {2023},
  note         = {Citation details to be verified}
}

@article{wangrosner2019bayesian,
  author       = {Wang, Chenguang and Rosner, Gary L.},
  title        = {A {B}ayesian nonparametric causal inference model for synthesizing randomized clinical trial and real-world evidence},
  year         = {2019},
  note         = {Citation details to be verified}
}

\end{filecontents*}

%==========================================================================
%  Document body
%==========================================================================

\title{Candidate additions to the introduction\\of \emph{Bayesian nonparametric fusion of randomised and real-world survival data for heterogeneous treatment effects}}
\author{Working note}
\date{\today}

\begin{document}
\maketitle

\noindent
This note inventories papers that could be added to the two literature paragraphs of the introduction (the confounding-function / frequentist paragraph and the Bayesian-borrowing paragraph). Organisation is by importance: the single closest competitor that is currently missing; references already in the working-paper bibliography that could be promoted to the introduction; references not yet on file that fit the narrative. BibTeX entries for every paper are embedded at the top of this file (the \verb|filecontents| block writes \verb|intro_additions.bib| on first compile).

\section*{1. The single most important omission}

\citet{yozova2025data} extend Bayesian Causal Forests to fuse trial and confounded observational data through an additively augmented BCF model with source-specific prognostic and treatment-effect components and a data-driven tempering of the observational likelihood. This is the closest Bayesian competitor to the proposed method: Bayesian, tree-based, fuses an RCT with a confounded observational source, and explicitly models the bias from unmeasured confounding. The differences are the outcome type (continuous rather than survival) and the bias-correction mechanism (tempered likelihood rather than a confounding-function decomposition under an AFT). A reasonable insertion in Paragraph~2 is: \emph{``\citet{yozova2025data} extend BCF to fuse trial and confounded observational data through a tempered-likelihood construction, but for continuous outcomes.''} The closing gap statement of Paragraph~2 sharpens once this paper is acknowledged explicitly.

\section*{2. Already in the working-paper references; candidates for promotion to the introduction}

\paragraph{Paragraph~1 (frequentist).}
\begin{itemize}
  \item \citet{chengcai2021adaptive} propose an adaptive combination of RCT and observational data for the ATE, with a data-driven weight that interpolates between trial-only and pooled estimators. Natural ATE precursor before the HTE-targeted Yang line.
  \item \citet{shyr2025multi} develop a multi-study R-learner for HTE estimation across multiple studies under within-study unconfoundedness. Belongs next to \citet*{wu2022integrative} and broadens the multi-source HTE strand.
  \item \citet{yangding2020combining} construct a unified estimating-equation framework for combining multiple observational datasets, optimally weighting sources to minimise asymptotic variance. A more general data-fusion precursor.
  \item \citet{vanderlaan2024adaptive} develop an adaptive targeted maximum-likelihood estimator that fuses RCT and RWD with a data-driven RWD-influence mechanism. Rounds out the frequentist ATE-fusion landscape.
  \item \citet{lintarpevans2024fusion} critically compare the confounding-function approach with experimental grounding under simulations for ATE estimation. Useful both as a comparative reference and because their simulations question whether either approach reliably beats RCT-only for ATE---which makes the HTE focus of the present work more pointed.
\end{itemize}

\paragraph{Paragraph~2 (Bayesian).}
\begin{itemize}
  \item \citet{muller2023bayesian} review Bayesian methodologies for integrating real-world data into clinical studies, covering power priors, MAP, commensurate priors and recent developments. Useful as the Bayesian counterpart of \citet{colnet2024causal}: a one-stop ``see X for review'' citation that avoids enumerating every variant.
  \item \citet{murray2015combining} extend commensurate priors to nonexchangeable functional and survival sources in oncology, using generalised mixture commensurate priors. Fits alongside \citet*{alt2025control} as a Bayesian-borrowing precursor for survival data.
  \item \citet{ohigashi2022horseshoe} and \citet{ohigashi2025potential} develop horseshoe and other continuous-shrinkage priors for dynamic borrowing of multiple historical controls. Optional additions to the ``later refinements'' lineage.
\end{itemize}

\section*{3. Already in the methodology references; candidates for promotion}

\begin{itemize}
  \item \citet{henderson2020individualized} develop a fully nonparametric Bayesian accelerated failure time model for individualised treatment effects with censored data, in a single source under unconfoundedness. This is the natural single-source Bayesian-AFT HTE precursor to the proposed method and is currently cited only in the methodology section. Promoting it to the introduction (last sentence of Paragraph~2, before the gap statement) tightens the gap argument: \emph{``\citet{henderson2020individualized} develop nonparametric Bayesian AFT modelling for individualised treatment effects in a single source.''}
  \item \citet{hahn2020bayesian} introduce Bayesian Causal Forests with a confounding-aware decomposition into prognostic and treatment-effect surfaces. Already cited in the methodology; a brief mention in the introduction would place the BCF lineage explicitly.
\end{itemize}

\section*{4. Not currently in the references; plausible additions}

\begin{itemize}
  \item \citet{hu2021estimating} compare Bayesian AFT-BART and other machine-learning methods for HTE estimation in observational survival data under unconfoundedness, finding that AFT-BART performs best on bias, precision and credible-interval coverage. The closest single-source Bayesian-survival-HTE precursor to the proposed method.
  \item \citet{zhao2023efficient} propose a doubly robust transfer-learning framework for the optimal individualised treatment regime under right censoring, combining trial and observational data with covariate shift. A related but distinct estimand (decision rule rather than CATE) and bias mechanism (covariate shift rather than unmeasured confounding); worth one sentence to delineate.
  \item \citet{colnet2024causal} provide a comprehensive review of frequentist data-fusion methods. Saves listing every variant on the frequentist side.
  \item \citet{bareinboim2016causal} establish the theory of identifiability for data-fusion problems. Useful at the start of Paragraph~1 if a foundational reference is desired.
  \item \citet{bi2023pamhc} construct a Bayesian nonparametric hybrid control (PAM-HC) for randomised clinical trials using external data. Pairs with \citet*{zhou2021incorporating} as a Bayesian-BART precursor.
  \item Suggested but with BibTeX details still to confirm: Lee et al. (2022, 2023, 2024) on doubly robust generalisability for survival; Alt et al.\ (2023) LEAP prior for dynamic Bayesian borrowing; Wang \& Rosner (2019) on Bayesian nonparametric multi-source propensity-score adjustment.
\end{itemize}

\section*{5. Test-then-pool and selective-borrowing strand}

A parallel strand to the confounding-function modelling addresses RWD bias by selection rather than by estimation. The canonical anchors are \citet*{viele2014use} for test-then-pool, \citet*{rosenman2023combining} for weighted combination, and \citet*{gaoyang2023adaptive} for adaptive selection. A single sentence in Paragraph~1 with two or three of these citations is sufficient to delineate the contrast with the confounding-function modelling approach taken in this work. (BibTeX entries to add once verified.)

\section*{Minimal recommendation}

If only three papers are added, the following are the highest-impact:
\begin{enumerate}
  \item \textbf{\citet{yozova2025data}} --- closest Bayesian competitor; currently missing.
  \item \textbf{\citet{henderson2020individualized}} --- natural single-source Bayesian-AFT HTE precursor; sharpens the final gap statement.
  \item \textbf{\citet{colnet2024causal}} or \textbf{\citet{lintarpevans2024fusion}} --- umbrella review citation, replaces enumeration on the frequentist side.
\end{enumerate}
A slightly expanded version adds \citet{shyr2025multi} to extend the HTE-fusion lineage and \citet{chengcai2021adaptive} to cover the adaptive-combination ATE precursor.

\bibliographystyle{plainnat}
\bibliography{intro_additions}

\end{document}